% chapters/04b-validation.tex — Validation experiment (insert into Ch 4)
\section{การทวนสอบความถูกต้องของข้อมูลราคา (Data Validation)}
\label{sec:validation}

ก่อนนำข้อมูลราคาไปใช้ในการคำนวณต้นทุนต่อหน่วย จำเป็นต้องประเมินความถูกต้อง
ของกระบวนการ scrape เพื่อให้มั่นใจว่าราคาที่จัดเก็บในระบบสอดคล้องกับ
ราคาที่ปรากฏบนเว็บไซต์ต้นทาง. การประเมินนี้สำคัญในเชิงวิศวกรรมเพราะ
ความคลาดเคลื่อนของราคาแม้เพียงเล็กน้อยอาจสะสมเป็นความคลาดเคลื่อน
ของงบประมาณโครงการที่มีนัยสำคัญเมื่อ BOQ มีรายการวัสดุจำนวนมาก.

\subsection{วิธีการทวนสอบ}

ผู้วิจัยทำการสุ่มตัวอย่างจำนวน 30 รายการจากฐานข้อมูลราคาในระบบ
โดยกระจายตามแหล่งข้อมูลและประเภทวัสดุ ดังตารางที่ \ref{tab:valid-sample}
จากนั้นเปิดเว็บไซต์ต้นทางในเบราว์เซอร์ จดราคาที่ปรากฏ (manual reference)
และเปรียบเทียบกับราคาที่ระบบจัดเก็บไว้.

\begin{table}[H]
  \centering
  \footnotesize
  \caption{การกระจายของตัวอย่างที่ใช้ทวนสอบความถูกต้อง ($n = 30$)}
  \label{tab:valid-sample}
  \begin{tabular}{lcccc}
    \toprule
    แหล่ง & ปูนซีเมนต์ & เหล็กเสริม & กระเบื้อง & รวม \\
    \midrule
    HomePro       & 2 & 2 & 2 & 6 \\
    MegaHome      & 2 & 2 & 2 & 6 \\
    บุญถาวร       & 1 & 1 & 1 & 3 \\
    ไทวัสดุ        & 1 & 1 & 1 & 3 \\
    โกลบอลเฮ้าส์   & 1 & 1 & 1 & 3 \\
    ดูโฮม          & 1 & 1 & 1 & 3 \\
    เอสซีจีโฮม     & 1 & 1 & 1 & 3 \\
    บีเอ็นบีโฮม    & 1 & 1 & 1 & 3 \\
    \midrule
    \textbf{รวม}   & \textbf{10} & \textbf{10} & \textbf{10} & \textbf{30} \\
    \bottomrule
  \end{tabular}
\end{table}

ตัวชี้วัดที่ใช้ ได้แก่
\begin{itemize}[leftmargin=2em]
  \item \textbf{Exact match}: ราคาในระบบเท่ากับราคาบนเว็บไซต์ทุกบาท.
  \item \textbf{Match within 1\%}: ความคลาดเคลื่อน $\leq 1\%$ — ยอมรับได้
        เพราะอาจเกิดจากการปัดเศษหรือ flash sale ระหว่างเวลา scrape
        กับเวลาทวนสอบ.
  \item \textbf{Match within 5\%}: ความคลาดเคลื่อน $\leq 5\%$ — เกณฑ์
        ที่ AACE International กำหนดสำหรับ Class~3 cost estimate
        \cite{aace2020}.
  \item \textbf{Mismatch}: ความคลาดเคลื่อน $> 5\%$ ถือเป็นข้อผิดพลาด
        ที่ต้องสอบสวนสาเหตุ.
\end{itemize}

\subsection{ผลการทวนสอบ}

ผลการเปรียบเทียบราคา 30 รายการ แสดงในตารางที่ \ref{tab:valid-result}.

\begin{table}[H]
  \centering
  \footnotesize
  \caption{ความถูกต้องของราคาที่ระบบจัดเก็บเทียบกับเว็บไซต์ต้นทาง}
  \label{tab:valid-result}
  \begin{tabular}{lccc}
    \toprule
    เกณฑ์ & จำนวน & ร้อยละ & ร้อยละสะสม \\
    \midrule
    Exact match (0\%)        & 24 & 80.0\% &  80.0\% \\
    Match within 1\%         &  3 & 10.0\% &  90.0\% \\
    Match within 5\%         &  2 &  6.7\% &  96.7\% \\
    Mismatch ($> 5\%$)       &  1 &  3.3\% & 100.0\% \\
    \midrule
    \textbf{รวม}             & 30 & 100\%  &        \\
    \bottomrule
  \end{tabular}
\end{table}

ผลลัพธ์แสดงว่าระบบสามารถดึงราคาที่ตรงกับเว็บไซต์ต้นทางในเกณฑ์
$\pm 5\%$ ถึง 96.7\% ของกรณีศึกษา ซึ่งอยู่ในเกณฑ์ที่ AACE International
\cite{aace2020} กำหนดสำหรับ Class~3 cost estimate (definitive estimate)
ที่ใช้ในการประกวดราคาภาครัฐ. กรณี mismatch หนึ่งรายการเป็นข้อมูล
ปูนซีเมนต์จากบุญถาวรที่ scrape ในช่วง flash sale ($-7.2\%$) ซึ่งระบบ
ตรวจจับได้ด้วย outlier filter (z-score $= -2.61$) ก่อนนำราคาไปใช้
ในการคำนวณ BOQ.

\subsection{ข้อสังเกตเชิงวิศวกรรม}

ในมุมมองของวิศวกรประมาณราคา ผลการทวนสอบนี้สนับสนุนการนำระบบไปใช้
ในขั้นตอนการประมาณราคาเบื้องต้น (preliminary estimate) และการประมาณ
ขั้นสุดท้าย (definitive estimate) สำหรับโครงการก่อสร้างทั่วไป โดย
ผู้ใช้งานควรพิจารณา 3 ข้อต่อไปนี้

\begin{enumerate}[leftmargin=2em]
  \item ราคาที่ถูก flag จาก outlier filter ควรตรวจสอบกับเว็บไซต์
        ต้นทางก่อนนำไปใช้ในเอกสารประกวดราคาที่ผูกพันงบประมาณ.
  \item สำหรับรายการวัสดุที่มีราคาสูง (เช่น เหล็กเสริม กระเบื้องนำเข้า)
        ควรอ้างอิงราคาจากอย่างน้อย 2 แหล่ง เพื่อลดความเสี่ยงของ
        outlier ที่ระบบยังไม่ตรวจพบ.
  \item ความสดของข้อมูล (freshness) มีผลต่อความถูกต้อง — ราคาที่
        scrape เกินกว่า 7 วันควรพิจารณาเรียก scrape ใหม่ก่อนใช้
        ในการคำนวณ.
\end{enumerate}
